%% ch05.tex --- Making Gradient Descent Work
%% Numbers are computed by verify/ch05_scaling.py.

\chapter{Making Gradient Descent Work}
\label{ch:ch05}

\epigraph{With four parameters I can fit an elephant, and with five I can make
him wiggle his trunk.}{---\ John von Neumann}

\begin{chapterlearning}
After working through this chapter you will be able to:
\begin{enumerate}
  \item Explain why features on different scales make gradient descent slow, in
        terms of the cost surface and its condition number.
  \item Apply mean normalization, min--max scaling and standardisation, and say
        which part of each does the useful work.
  \item Predict the effect of scaling on the number of iterations required, and
        state what scaling does \emph{not} change.
  \item Choose a learning rate with the $\times 3$ ladder and read the result.
  \item Diagnose a training run from the shape of its learning curve.
  \item State the rule for computing scaling statistics, and say exactly when
        breaking it changes the model.
  \item Construct engineered features from domain knowledge, and explain why a
        model linear in its parameters can fit a curve.
  \item Fit a polynomial regression, and recognise the point at which added
        flexibility stops helping.
\end{enumerate}
\end{chapterlearning}

\Cref{ch:ch04} ended with a loose end. Its worked example fitted five apartments
by gradient descent, but only after standardising the two features, with a
parenthetical promise that this chapter would explain why. The explanation turns
out to be worth a chapter, because the same idea --- that the \emph{shape} of the
cost surface, not just its minimum, determines what happens --- runs through
everything that follows.

\section{The Problem, Measured}
\label{sec:scaling-problem}

Dataset~D2 has two features on comparable-looking ranges: area from $8$ to $22$,
age from $5$ to $20$. Nothing about them looks troublesome. Yet the Hessian of
its cost has eigenvalues
\[
  \lambda_{\min} = 0.01315,
  \qquad
  \lambda_{\max} = 352.2,
\]
a ratio of $26{,}775$. By \cref{eq:alpha-star} the largest usable learning rate
is $\alpha^\star = 2/\lambda_{\max} = 0.005678$, and by the reasoning of
\cref{sec:learning-rate} the slowest component of the error decays by a factor
of only $1 - \alpha\lambda_{\min} \approx 0.99993$ per iteration. That is the
whole difficulty in two numbers: steps must be tiny because of the largest
eigenvalue, and progress is glacial because of the smallest.

\begin{defbox}{Condition number}
The \term{condition number} of the cost is
$\kappa = \lambda_{\max}/\lambda_{\min}$, the ratio of the largest to the
smallest eigenvalue of its Hessian. It is the axis ratio of the elliptical
contours, squared\footnote{The contours are ellipses whose axis lengths go as
$1/\sqrt{\lambda}$, so the axis ratio is $\sqrt{\kappa}$.}, and the number of
gradient descent iterations required grows roughly in proportion to it.
\end{defbox}

\section{Feature Scaling}
\label{sec:feature-scaling}

The cure is to change the units of the features. Three transformations are in
common use, all applied one feature (one column of $\Xmat$) at a time.

\begin{defbox}{Three scalings}
For feature $j$ with values $\xjj{1}{j},\dots,\xjj{m}{j}$, mean $\mu_j$,
standard deviation $\sigma_j$, minimum $\ell_j$ and maximum $h_j$:
\[
  \underbrace{x_j \mapsto \frac{x_j}{h_j}}_{\text{divide by the maximum}},
  \qquad
  \underbrace{x_j \mapsto \frac{x_j - \mu_j}{h_j - \ell_j}}_{\text{mean normalization}},
  \qquad
  \underbrace{x_j \mapsto \frac{x_j - \mu_j}{\sigma_j}}_{\text{standardisation, or $z$-scoring}} .
\]
The first two put every feature on a range of roughly unit width; the last gives
every feature mean $0$ and variance $1$. All three are affine, so all three are
undone by the model's own parameters --- a point \cref{sec:scaling-payoff}
makes precise.
\end{defbox}

Applied to D2, they do markedly different amounts of good:

\begin{center}
\small
\renewcommand{\arraystretch}{1.2}
\begin{tabular}{@{}lrrrr@{}}
\toprule
scaling & $\lambda_{\min}$ & $\lambda_{\max}$ & $\kappa$ & $\alpha^\star$\\
\midrule
raw                & $0.01315$ & $352.21$ & $26{,}775.0$ & $0.005678$\\
divide by maximum  & $0.00579$ & $1.7719$ & $306.1$      & $1.128759$\\
mean normalized    & $0.02176$ & $1.0000$ & $46.0$       & $2.000000$\\
standardised       & $0.18754$ & $1.8125$ & $9.7$        & $1.103474$\\
\bottomrule
\end{tabular}
\end{center}

Read the second row carefully, because it contradicts the advice usually given.
Dividing by the maximum puts both features on $[0,1]$ --- it does exactly what
``get your features onto similar ranges'' asks --- and yet leaves the condition
number at $306$, thirty times worse than standardisation. Matching the ranges is
not the operation that matters most.

\subsection{What actually does the work}

What matters most is \emph{centring}. The design matrix carries a column of ones
for the bias, and if a feature never comes near zero, that column and the feature
column point in nearly the same direction: they are close to linearly dependent,
which is precisely the near-singularity of \cref{sec:multicollinearity}. Area
runs from $8$ to $22$, so the area column is roughly $15$ times the ones column
plus a small variation --- and $\lambda_{\min} = 0.013$ is the result.

Subtracting the mean removes the dependency exactly. Once every feature has mean
zero, each feature column is orthogonal to the ones column, the Hessian becomes
block diagonal, and the bias direction decouples from the rest.

\begin{derivation}{What is left after centring}
Centre and standardise every feature. Then $\Xmat^{\mathsf{T}}\mathbf{1} =
\mathbf{0}$ and each feature has unit variance, so
\[
  H = \frac1m \begin{pmatrix} m & \mathbf{0}^{\mathsf{T}}\\
                              \mathbf{0} & \Xmat^{\mathsf{T}}\Xmat\end{pmatrix}
    = \begin{pmatrix} 1 & \mathbf{0}^{\mathsf{T}}\\
                      \mathbf{0} & R\end{pmatrix},
\]
where $R$ is the \emph{correlation matrix} of the features. For two features
with correlation $\rho$, $R = \left(\begin{smallmatrix}1&\rho\\ \rho&1
\end{smallmatrix}\right)$, whose eigenvalues are $1+|\rho|$ and $1-|\rho|$.

For D2 the two features have $\rho = -0.8125$: in this small sample the larger
apartments happen also to be the newer ones. Hence the eigenvalues $1.8125$ and
$0.1875$, and $\kappa = 9.67$, exactly the last row of the table.
\end{derivation}

That derivation also marks the limit of what scaling can achieve. After
standardisation the condition number is set entirely by the correlations between
features, and no amount of rescaling individual columns will change a
correlation. Removing that last factor of ten requires either decorrelating the
features (\cref{ch:ch16}) or adding a penalty that bounds the parameters
(\cref{ch:ch07}).

\subsection{The picture}

Two features and a bias give a three-dimensional parameter space that cannot be
drawn, so return to the four-point set of \cref{ch:ch02}, where there is only
$(w,b)$. Its feature is $x = 1,2,3,4$, with mean $2.5$ and standard deviation
$1.1180$.

\begin{figure}[htbp]
  \centering
  \begin{minipage}[t]{0.48\linewidth}
    \centering
    \begin{tikzpicture}
      \begin{axis}[width=\linewidth, height=4.3cm,
        xlabel={$w$}, ylabel={$b$},
        tick label style={font=\scriptsize}, label style={font=\small},
        view={0}{90}, domain=0.2:3.0, y domain=-3:2,
        colormap={ml}{color=(MLBlue) color=(MLTeal) color=(MLAmber)}]
        \costcontour{}{0.1,0.25,0.6,1.2,2.2,3.6}
          {(30*x^2 + 20*x*y + 4*y^2 - 86*x - 28*y + 62)/8}
        \addplot3[MLRose,thick,mark=*,mark size=0.9pt] coordinates {
          (0.000,0.000,0)(1.287,0.419,0)(1.293,0.403,0)(1.298,0.386,0)
          (1.304,0.371,0)(1.309,0.355,0)(1.314,0.340,0)(1.325,0.310,0)
          (1.334,0.281,0)(1.348,0.240,0)(1.362,0.201,0)(1.378,0.152,0)
          (1.397,0.095,0)(1.415,0.044,0)(1.446,-0.046,0)};
        \addplot3[only marks,mark=star,mark size=2.6pt,MLBlue]
          coordinates {(1.6,-0.5,0)};
      \end{axis}
    \end{tikzpicture}
    {\scriptsize Raw: $\kappa = 55.8$, \textbf{222 iterations}.\par}
  \end{minipage}\hfill
  \begin{minipage}[t]{0.48\linewidth}
    \centering
    \begin{tikzpicture}
      \begin{axis}[width=\linewidth, height=4.3cm,
        xlabel={$w$}, ylabel={$b$},
        tick label style={font=\scriptsize}, label style={font=\small},
        view={0}{90}, domain=-0.6:4.2, y domain=1.1:5.9,
        colormap={ml}{color=(MLBlue) color=(MLTeal) color=(MLAmber)}]
        \costcontour{}{0.1,0.25,0.6,1.2,2.2,3.6}
          {(4*x^2 + 4*y^2 - 14.3108*x - 28*y + 62)/8}
        \addplot3[MLRose,thick,mark=*,mark size=0.9pt] coordinates {
          (0.000,0.000,0)(1.789,3.500,0)};
        \addplot3[only marks,mark=star,mark size=2.6pt,MLBlue]
          coordinates {(1.789,3.5,0)};
      \end{axis}
    \end{tikzpicture}
    {\scriptsize Standardised: $\kappa = 1.0$, \textbf{1 iteration}.\par}
  \end{minipage}
  \caption{The same four points and the same step-size rule
  ($\alpha = 1/\lambda_{\max}$). Left: elliptical contours, and a descent that
  crawls along the valley floor. Right: standardising $x$ makes the contours
  exact circles, so the negative gradient points straight at the minimum and one
  step arrives. The optimum has \emph{moved}, from $(1.6,-0.5)$ to
  $(1.7889,\,3.5)$, because the parameters now multiply different units.}
  \label{fig:scaling-contours}
\end{figure}

The circular case is not a coincidence of this dataset. With one standardised
feature, $\sum_i z^{(i)} = 0$ and $\frac1m\sum_i (z^{(i)})^2 = 1$, so
$H = I$ exactly, whatever the data was. Gradient descent on a perfectly circular
bowl with $\alpha = 1$ lands on the minimum in a single step from any starting
point.

\section{The Payoff, and What It Is Not}
\label{sec:scaling-payoff}

Back to D2, and the question \cref{ch:ch04} left open. Running gradient descent
under each scaling at $90\%$ of that scaling's own stability limit, and counting
iterations until the cost is within $1\%$ of its minimum:

\begin{center}
\small
\renewcommand{\arraystretch}{1.2}
\begin{tabular}{@{}lrr@{}}
\toprule
scaling & $\alpha$ used & iterations\\
\midrule
raw               & $0.005111$ & $26{,}605$\\
divide by maximum & $1.015883$ & $521$\\
mean normalized   & $1.800000$ & $89$\\
standardised      & $0.993126$ & $23$\\
\bottomrule
\end{tabular}
\end{center}

Twenty-three iterations against twenty-six thousand: a factor of $1150$, from
three lines of arithmetic applied before training starts. That is the answer to
\cref{ch:ch04}'s question. Its worked example did not standardise because the
answer would otherwise be wrong; it standardised because the run would otherwise
have taken a thousand times longer to become right.

\begin{keyinsight}{Scaling changes the optimization, not the model}
Every scaling in this chapter is affine, and the model can absorb any affine
change of variables into its own parameters: if $z = (x-\mu)/\sigma$ then
$w_z z + b_z$ is the same function of $x$ as $w x + b$ with
$w = w_z/\sigma$ and $b = b_z - w_z\mu/\sigma$. Fitted to convergence, an
unregularized least-squares model gives \emph{identical predictions} whether you
scale or not --- in the experiment of \cref{sec:scaling-practice}, identical to
$1.4\times 10^{-14}$.

So scaling is not a modelling improvement, and a model that fits badly will not
fit better for being scaled. It is an improvement to the \emph{path} the
optimizer takes to the same destination. The exceptions --- the cases where
scaling does change the answer --- are exactly the ones in
\cref{sec:scaling-practice}, and they are worth knowing precisely because the
general rule is otherwise so clean.
\end{keyinsight}

Note also that the fitted parameters do change, as \cref{fig:scaling-contours}
shows: $w$ moved from $1.6$ to $1.7889$. It must, because it now measures the
change in $y$ per standard deviation of $x$ rather than per unit of $x$. If you
scale, you must either report coefficients in scaled units and say so, or map
them back with the formulas above.

\section{Choosing the Learning Rate}
\label{sec:alpha-ladder}

\Cref{sec:learning-rate} gave the theory: $\alpha < 2/\lambda_{\max}$. In
practice nobody computes $\lambda_{\max}$, and the working procedure is a
search over a coarse grid.

\begin{algorithmbox}{Standardise, then find $\alpha$ by the $\times 3$ ladder}
\begin{enumerate}\setlength{\itemsep}{1pt}
  \item For each feature $j$, compute $\mu_j$ and $\sigma_j$ \emph{on the
        training set only}, and store them.
  \item Replace each $\xjj{i}{j}$ by $(\xjj{i}{j}-\mu_j)/\sigma_j$.
  \item For $\alpha \in \{0.001,\,0.003,\,0.01,\,0.03,\,0.1,\,0.3,\,1\}$: run a
        few hundred iterations and record the cost curve.
  \item Discard every $\alpha$ whose cost rises or oscillates upward.
  \item Among the rest, take the largest --- or, cautiously, the one below it.
  \item Train properly at that $\alpha$, applying the stored $\mu_j,\sigma_j$ to
        every future input.
\end{enumerate}
\end{algorithmbox}

Two hundred iterations of that ladder on D2, before and after standardisation:

\begin{center}
\small
\renewcommand{\arraystretch}{1.15}
\begin{tabular}{@{}lr@{\qquad}lr@{}}
\toprule
\multicolumn{2}{c}{\textbf{raw features}} &
\multicolumn{2}{c}{\textbf{standardised}}\\
$\alpha$ & $J$ after $200$ & $\alpha$ & $J$ after $200$\\
\midrule
$0.00001$ & $518.10$   & $0.003$ & $524.306$\\
$0.00003$ & $96.62$    & $0.01$  & $35.956$\\
$0.0001$  & $22.11$    & $0.03$  & $2.0916$\\
$0.0003$  & $1.8275$   & $0.1$   & $0.894917$\\
$0.001$   & $1.2054$   & $0.3$   & $\mathbf{0.889051}$\\
$0.003$   & $1.2021$   & $1.0$   & $\mathbf{0.889051}$\\
$0.01$    & diverged   & $3.0$   & diverged\\
\bottomrule
\end{tabular}
\end{center}

The two columns tell different stories. On the right, $\alpha = 0.3$ reaches the
exact minimum $0.889051$ within two hundred iterations, and the ladder has
plainly found it. On the left, the best surviving value, $\alpha = 0.003$, still
reports $1.2021$ --- thirty-five per cent above the minimum --- and the next rung
up diverges. There is no good rung. A ladder run on unscaled features can leave
you concluding that the model fits poorly when in fact the optimizer has not
finished.

\section{Reading the Learning Curve}
\label{sec:learning-curve}

\Cref{sec:stopping} covered when to stop. The complementary skill is diagnosing
what went wrong from the shape of $J$ against iteration, plotted --- as that
section insisted --- on a logarithmic vertical axis.

\begin{center}
\small
\renewcommand{\arraystretch}{1.3}
\begin{tabular}{@{}L{4.0cm}L{7.0cm}@{}}
\toprule
\textbf{What you see} & \textbf{What it means}\\
\midrule
Smooth decrease, still falling at the last iteration &
  Working; $\alpha$ may be too small, or the budget too short. Try the next rung
  up, or more iterations.\\
Smooth decrease, flat for the last third &
  Converged, or stuck on a valley floor. Check the gradient norm, not the cost
  (\cref{sec:stopping}).\\
Decrease with visible oscillation &
  $\alpha$ near the stability limit. Usually still converging, sometimes fastest
  of all; drop one rung if it worries you.\\
Increase, or oscillation that grows &
  $\alpha > \alpha^\star$. Drop two rungs.\\
Immediate \code{nan} or \code{inf} &
  Divergence to overflow, or unscaled features with a large $\alpha$, or a bug
  in the gradient. Check the gradient against a finite-difference estimate.\\
Falls, then rises steadily &
  Not an optimization problem at all if the rising curve is the \emph{validation}
  cost --- that is overfitting (\cref{ch:ch07}).\\
\bottomrule
\end{tabular}
\end{center}

\section{Scaling in Practice}
\label{sec:scaling-practice}

The rule is short: \textbf{compute $\mu_j$ and $\sigma_j$ on the training set,
and apply those same numbers to validation data, to test data, and to every
input the deployed model ever sees.} Never recompute them on the data you are
predicting.

The first reason is practical and unanswerable: at prediction time there may be
only one input, which has no standard deviation. Whatever you scale by must have
been decided in advance.

The second reason is the one usually given --- that using test-set statistics
leaks information from the test set into training. That is true, but it is worth
being precise about how much damage it does, because the honest answer is
sometimes none at all.

\begin{workedexample}{When does the choice of statistics change the model?}
Two features, sixty training examples, twenty test examples. Fit twice: once
scaling by training statistics, once by statistics computed over all eighty
examples.

\begin{center}
\small
\renewcommand{\arraystretch}{1.2}
\begin{tabular}{@{}lrrr@{}}
\toprule
model & train stats & all-data stats & max $|\text{difference}|$\\
\midrule
least squares, converged & $1.656872$ & $1.656872$ & $1.4\times10^{-14}$\\
ridge, $\lambda = 1$     & $1.584349$ & $1.580819$ & $0.059$\\
ridge, $\lambda = 10$    & $2.852947$ & $3.064429$ & $0.429$\\
gradient descent, $20$ iterations & $1.631928$ & $1.615459$ & $0.097$\\
gradient descent, $200$ iterations & $1.656872$ & $1.656872$ & $0.000$\\
\bottomrule
\end{tabular}
\end{center}
(Test RMSE in the first two columns; the last column is the largest difference
between the two vectors of predictions.)

For ordinary least squares fitted to convergence, the two are the same model to
machine precision --- as the Key Insight above promised, since the scaling is
affine and the fit absorbs it. Add a penalty that depends on the size of the
coefficients (\cref{ch:ch07}), and the two diverge, because the penalty is
measured in scaled units and the units now depend on the test set. Stop the
optimizer early instead of running it to convergence, and they diverge again,
for the same reason: the path depends on the scaling even though the destination
does not.
\end{workedexample}

\begin{pitfall}{Fit the scaler on the training set, always}
The exception above is a trap, not a licence. Almost every model you will
actually deploy is regularized, or early-stopped, or uses distances
(\cref{ch:ch13}), or feeds a pipeline where something further downstream is
scale-dependent. In all of those, statistics computed over the full dataset make
the test score optimistic and the deployed model different from the one you
evaluated.

The rule costs nothing to follow and removes an entire category of
irreproducible result, so follow it unconditionally. In \tool{scikit-learn} this
means \code{scaler.fit(X\_train)} followed by \code{scaler.transform(X\_test)},
never \code{fit\_transform} on both --- or, better, a \code{Pipeline}, which
makes the mistake impossible to express.
\end{pitfall}

\section{Feature Engineering}
\label{sec:feature-engineering}

So far the features have been whatever the data happened to contain. Often the
most valuable thing you can do to a model is give it a feature nobody recorded.

\begin{workedexample}{Frontage and depth, or area?}
Six building plots are recorded by frontage and depth, and sold for a price:

\begin{center}
\small
\renewcommand{\arraystretch}{1.15}
\begin{tabular}{@{}rrrr@{}}
\toprule
frontage & depth & area $=$ f\,$\times$\,d & price\\
\midrule
$20$ & $50$ & $1000$ & $51$\\
$30$ & $40$ & $1200$ & $59$\\
$25$ & $60$ & $1500$ & $76$\\
$40$ & $30$ & $1200$ & $61$\\
$35$ & $45$ & $1575$ & $78$\\
$45$ & $50$ & $2250$ & $113$\\
\bottomrule
\end{tabular}
\end{center}

A linear model in frontage and depth cannot express a product, and it shows:

\begin{center}
\small
\begin{tabular}{@{}lr@{}}
\toprule
features & $J$\\
\midrule
frontage, depth        & $2.5190$\\
area alone             & $0.3576$\\
frontage, depth, area  & $0.1523$\\
\bottomrule
\end{tabular}
\end{center}

One engineered feature, replacing two recorded ones, cuts the cost by a factor
of seven --- and does it with \emph{fewer} parameters. No change of algorithm,
no extra data, no tuning: just the observation that land is sold by the square
foot, which any estate agent could have told you.
\end{workedexample}

Engineered features that repay the effort, in rough order of frequency:

\begin{itemize}
  \item \textbf{Products}, when two quantities act together: frontage $\times$
        depth, price $\times$ quantity, dose $\times$ duration.
  \item \textbf{Ratios}, when the meaningful quantity is relative: price per
        square foot, errors per thousand requests, debt to income.
  \item \textbf{Differences}, especially of dates: age at purchase rather than
        date of birth and date of purchase.
  \item \textbf{Aggregates over a history}: a customer's mean order value, the
        number of failed logins in the past hour.
  \item \textbf{Decompositions of a timestamp}: hour of day, day of week, whether
        it is a holiday. A raw timestamp is nearly useless as a number; its parts
        are informative.
\end{itemize}

\begin{keyinsight}{Linear in the parameters, not in the features}
``Linear regression'' constrains how the model depends on $\vw$ and $b$, not on
the data. Nothing stops a feature from being $x_1x_2$, or $\sqrt{x}$, or
$\log(1+x)$, or the day of the week extracted from a date. The model remains
linear in the parameters, so every derivation in \cref{ch:ch02,ch:ch03,ch:ch04}
--- the convexity, the closed form, the gradient --- continues to hold exactly.
This one observation is what makes the humble linear model as widely useful as it
is, and it is the bridge to the next section.
\end{keyinsight}

\section{Polynomial Regression}
\label{sec:polynomial}

The same trick fits curves. Given a single feature $x$, construct
$x, x^2, x^3, \dots$ as separate features and fit a linear model to them:
\[
  f(x) \;=\; w_1x + w_2x^2 + w_3x^3 + b .
\]
This is a cubic in $x$ and a linear model in $(w_1,w_2,w_3,b)$, so it is fitted
by exactly the machinery already built.

\begin{workedexample}{Ten points on a curve}
Dataset~T5, ten points on a curve:

\begin{center}
\small
\renewcommand{\arraystretch}{1.15}
\begin{tabular}{@{}l*{10}{r}@{}}
\toprule
$x$ & $1$ & $2$ & $3$ & $4$ & $5$ & $6$ & $7$ & $8$ & $9$ & $10$\\
$y$ & $2.1$ & $3.4$ & $5.6$ & $9.1$ & $13.2$ & $19.4$ & $25.9$ & $34.2$
    & $43.6$ & $54.5$\\
\bottomrule
\end{tabular}
\end{center}

\begin{center}
\small
\renewcommand{\arraystretch}{1.2}
\begin{tabular}{@{}lrr@{}}
\toprule
model & $J$ & $\kappa(\Xmat^{\mathsf{T}}\Xmat)$\\
\midrule
degree $1$ (a straight line) & $9.523067$ & $1.9\times10^{2}$\\
degree $2$                   & $0.019067$ & $4.1\times10^{4}$\\
degree $3$                   & $0.008885$ & $1.2\times10^{7}$\\
degree $4$                   & $0.008885$ & $4.4\times10^{9}$\\
degree $9$                   & $0.000000$ & $7.3\times10^{21}$\\
\bottomrule
\end{tabular}
\end{center}

The straight line is hopeless and the quadratic is excellent: adding $x^2$ drops
the cost by a factor of $500$. The cubic buys a further factor of two, the
quartic buys nothing at all, and the ninth-degree fit achieves exactly zero.
\end{workedexample}

Two things in that table deserve attention, and they point in opposite
directions.

\subsection{Powers make scaling compulsory}

The condition number rises by roughly two orders of magnitude per degree, and the
reason is visible in the feature ranges:

\begin{center}
\small
\begin{tabular}{@{}lrrr@{}}
\toprule
feature & minimum & maximum & ratio\\
\midrule
$x$   & $1$ & $10$      & $10$\\
$x^2$ & $1$ & $100$     & $100$\\
$x^3$ & $1$ & $1000$    & $1000$\\
$x^4$ & $1$ & $10{,}000$ & $10{,}000$\\
\bottomrule
\end{tabular}
\end{center}

A model with $x$ and $x^3$ as features has one input ranging over a decade and
another over three, which is the situation of \cref{sec:scaling-problem} made
deliberately worse. Standardisation is not optional here; without it, gradient
descent on a cubic fit is unusable, and by degree $9$ the condition number of
$10^{21}$ exceeds what double precision can represent faithfully, so even the
direct solver is returning something it cannot vouch for.

\subsection{Where the flexibility goes wrong}

The degree-$9$ fit passes exactly through all ten points, so $J = 0$. By the cost
function --- the only thing the algorithm can see --- it is the best model in the
table by an infinite margin. \Cref{fig:polynomial} shows what it actually is.

\begin{figure}[htbp]
  \centering
  \begin{tikzpicture}
    \begin{axis}[mlaxis, width=0.78\linewidth, height=5.6cm,
      xlabel={$x$}, ylabel={$y$},
      xmin=0.4, xmax=10.8, ymin=-22, ymax=95,
      legend style={font=\tiny, at={(0.02,0.98)}, anchor=north west}]
      \addplot[MLAmber,thick,dashed] coordinates {
        (0.60,-7.189)(1.60,-1.416)(2.60,4.357)(3.60,10.131)(4.60,15.904)
        (5.60,21.677)(6.60,27.451)(7.60,33.224)(8.60,38.997)(9.60,44.771)
        (10.60,50.544)};
      \addlegendentry{degree 1, $J=9.52$}
      \addplot[MLBlue,very thick] coordinates {
        (0.60,2.267)(1.10,2.363)(1.60,2.760)(2.10,3.457)(2.60,4.453)
        (3.10,5.750)(3.60,7.347)(4.10,9.243)(4.60,11.440)(5.10,13.937)
        (5.60,16.733)(6.10,19.830)(6.60,23.227)(7.10,26.923)(7.60,30.920)
        (8.10,35.217)(8.60,39.813)(9.10,44.710)(9.60,49.907)(10.10,55.403)
        (10.60,61.200)};
      \addlegendentry{degree 2, $J=0.019$}
      \addplot[MLRose,thick] coordinates {
        (0.60,-19.716)(0.85,-2.623)(1.10,3.814)(1.35,5.141)(1.60,4.558)
        (1.85,3.716)(2.10,3.306)(2.35,3.476)(2.60,4.110)(2.85,5.010)
        (3.10,5.994)(3.35,6.944)(3.60,7.814)(3.85,8.623)(4.10,9.424)
        (4.35,10.289)(4.60,11.277)(4.85,12.428)(5.10,13.747)(5.35,15.215)
        (5.60,16.789)(5.85,18.417)(6.10,20.054)(6.35,21.673)(6.60,23.277)
        (6.85,24.897)(7.10,26.591)(7.35,28.427)(7.60,30.467)(7.85,32.737)
        (8.10,35.207)(8.35,37.771)(8.60,40.261)(8.85,42.472)(9.10,44.263)
        (9.35,45.710)(9.60,47.366)(9.85,50.650)(10.10,58.394)(10.35,75.602)
        (10.60,110.445)};
      \addlegendentry{degree 9, $J=0$}
      \addplot[only marks,mark=*,mark size=1.6pt,MLGray,forget plot]
        coordinates {(1,2.1)(2,3.4)(3,5.6)(4,9.1)(5,13.2)(6,19.4)(7,25.9)
                     (8,34.2)(9,43.6)(10,54.5)};
    \end{axis}
  \end{tikzpicture}
  \caption{Three fits to the same ten points. The degree-$9$ curve has zero
  training cost --- it passes exactly through every point --- and is plainly the
  worst description of the data. Between $x=1$ and $x=2$ it dips and rises where
  the data does neither, and just outside the range of the data it leaves
  altogether.}
  \label{fig:polynomial}
\end{figure}

Asked to predict slightly outside the observed range, the two models part
company completely:

\begin{center}
\small
\begin{tabular}{@{}rrr@{}}
\toprule
$x$ & degree $2$ & degree $9$\\
\midrule
$1.5$  & $2.657$  & $4.887$\\
$5.5$  & $16.150$ & $16.150$\\
$9.5$  & $48.843$ & $46.613$\\
$10.5$ & $60.017$ & $93.699$\\
\bottomrule
\end{tabular}
\end{center}

At $x = 10.5$, half a unit beyond the last observation, the degree-$9$ model
predicts $93.7$ against the quadratic's $60.0$. Its largest coefficient has
magnitude $506$, on data whose values never exceed $55$: enormous coefficients
that very nearly cancel, which is the signature we first met in
\cref{sec:multicollinearity} and will meet again, named, in \cref{ch:ch07}.

\begin{pitfall}{Zero training cost is a warning, not an achievement}
Every model in the table above was fitted correctly, by a method that provably
finds the global minimum of the cost it was given. The cost was minimised; the
model still got worse. The training cost measures agreement with data the model
has already seen, and past some point the only way to improve it is to
contort the model into shapes the data does not support.

The remedy is the one \cref{ch:ch01} insisted on and this chapter has quietly
been ignoring: evaluate on data the model was not fitted to. \Cref{ch:ch07} makes
that systematic, gives the phenomenon its name --- overfitting --- and supplies
the standard defence.
\end{pitfall}

Powers are not the only useful transformation, and often not the best one. If a
quantity grows quickly, $\log(1+x)$ compresses it; if it saturates, $\sqrt{x}$
often fits better than a quadratic and cannot turn back downward; if it is
periodic, $\sin$ and $\cos$ of the phase are the natural features. Choose the
transformation from what you know about the quantity, not by increasing the
degree until the cost stops falling.

\begin{labbox}{Lab 5 --- Scaling, learning rates, and curves}
Implement \code{zscore(X)} returning the scaled matrix together with $\mu$ and
$\sigma$, and a \code{transform} that applies stored statistics to new data.
Reproduce the conditioning table of \cref{sec:feature-scaling} for D2 --- your
raw condition number must be $26{,}775$ and your standardised one $9.7$ --- and
the iteration counts of \cref{sec:scaling-payoff}. Run the $\times 3$ ladder on
both versions and confirm that no rung works well on the raw features. Then fit
polynomials of degree $1$ to $9$ to T5, on standardised powers, and plot all
three curves of \cref{fig:polynomial}. Finally, hold out the last two points,
refit each degree on the first eight, and report the error on the two held out:
this is \cref{ch:ch10} arriving early, and the ordering of the results will
surprise you if the rest of the chapter has not.
\end{labbox}

\section*{Chapter Summary}
\addcontentsline{toc}{section}{Chapter Summary}

\begin{itemize}
  \item The condition number $\kappa = \lambda_{\max}/\lambda_{\min}$ of the cost
        governs how many iterations gradient descent needs. For raw D2 it is
        $26{,}775$; standardised, $9.7$.
  \item Centring matters more than range-matching. Dividing by the maximum puts
        features on $[0,1]$ and still leaves $\kappa = 306$, because an uncentred
        feature column is nearly parallel to the bias column.
  \item After standardisation the Hessian is the correlation matrix of the
        features bordered by a $1$, so what remains is set by correlations
        between features --- which scaling cannot fix, and
        \cref{ch:ch07,ch:ch16} can.
  \item The payoff is large: $26{,}605$ iterations against $23$ on D2. With a
        single standardised feature the contours are exact circles and one step
        of length $\alpha = 1$ reaches the minimum.
  \item Scaling changes the optimization, not the model. Fitted to convergence
        without regularization, scaled and unscaled models predict identically to
        machine precision. Fitted parameters do change, and must be reported in
        the units they were fitted in.
  \item Choose $\alpha$ with a $\times 3$ ladder, discarding any rung whose cost
        rises and taking the largest that survives. On badly scaled data no rung
        may be satisfactory, which is itself the diagnosis.
  \item Compute scaling statistics on the training set and apply them everywhere
        else. Breaking this changes the model whenever a penalty, an early stop,
        or a distance is involved --- which is nearly always.
  \item Linear regression is linear in its parameters, not its features. Products,
        ratios, differences and powers are all admissible features, and every
        result derived so far survives.
  \item Polynomial features make scaling compulsory, since $x^k$ ranges over $k$
        decades. Enough degrees will drive the training cost to zero while making
        the model worse --- the subject of \cref{ch:ch07}.
\end{itemize}

\section*{Exercises}
\addcontentsline{toc}{section}{Exercises}

\exercisesection{Short Questions}

\sq Define the condition number of a cost function and say what it predicts.

\sq Write the formulas for mean normalization and for standardisation, and give
the range or the moments each produces.

\sq Why does dividing every feature by its maximum leave the condition number
much larger than standardisation does?

\sq Explain why an uncentred feature column is nearly parallel to the bias
column, and what that does to $\lambda_{\min}$.

\sq After standardisation, what determines the remaining condition number, and
what can be done about it?

\sq Does feature scaling change the predictions of a converged, unregularized
linear model? Does it change the fitted parameters?

\sq State the $\times 3$ ladder procedure in three sentences.

\sq Give the rule for computing scaling statistics, and one practical reason it
cannot be broken even in principle.

\sq Name three kinds of engineered feature and give an example of each.

\sq In what sense is polynomial regression a linear model?

\exercisesection{Analytical Problems}

\ap A single feature is standardised to $z = (x-\mu)/\sigma$, and a model
$f(z) = w_z z + b_z$ is fitted. Derive the parameters $w, b$ of the equivalent
model $f(x) = wx + b$, and verify your formulas against the four-point set, where
$(w_z, b_z) = (1.7889, 3.5)$, $\mu = 2.5$ and $\sigma = 1.1180$ correspond to
$(w,b) = (1.6,-0.5)$.

\ap Show that after centring every feature, $\Xmat^{\mathsf{T}}\mathbf{1} =
\mathbf{0}$, and hence that the Hessian of the cost is block diagonal with the
bias in a block of its own. What is that block's single eigenvalue?

\ap For two standardised features with correlation $\rho$, show that the
eigenvalues of the correlation matrix are $1\pm|\rho|$, and hence that
$\kappa = (1+|\rho|)/(1-|\rho|)$. At what value of $|\rho|$ does $\kappa$ reach
$100$?

\ap A feature is measured in metres and standardised. The data is then
re-collected in centimetres and standardised again. Show that the standardised
values are identical, and explain why this makes standardisation preferable to
dividing by the maximum when units are arbitrary.

\ap The four-point set has $x = 1,2,3,4$. Compute $\mu$, $\sigma$, the
standardised values, and by hand verify that the resulting Hessian is the
identity. Then show that gradient descent with $\alpha=1$ from any starting
point reaches the minimum in one step.

\ap A degree-$d$ polynomial is fitted to $m$ points. For what $d$ can the
training cost be driven to exactly zero, and how many such fits are there when
$d$ exceeds that value?

\exercisesection{Design Problems}

\dpr You are handed a dataset whose features include an age in years, an income
in local currency, a binary indicator, a count of prior visits, and a latitude.
Decide how to scale each, justify any you would leave alone, and say what you
would store to make the transformation reproducible at prediction time.

\dpr A colleague reports that standardising their features improved test accuracy
by four per cent. Given this chapter, list the possible explanations in order of
likelihood, and describe the single experiment that would distinguish between
them.

\dpr You are modelling electricity demand from a timestamp, an outdoor
temperature and a holiday flag. Propose six engineered features, say what
knowledge each encodes, and identify which pairs you would expect to be strongly
correlated.

\dpr Write the checklist you would apply before every gradient-descent training
run --- scaling, learning rate, budget, stopping rule, diagnostics to plot ---
and for each item state what evidence would tell you it was wrong.

\exercisesection{Programming Exercises}

\pe Implement standardisation and reproduce the conditioning table of
\cref{sec:feature-scaling} for D2, including all four rows.

\pe Reproduce the iteration counts of \cref{sec:scaling-payoff}. Then plot
iterations required against condition number for a family of synthetic datasets
with $\kappa$ ranging over four orders of magnitude, on log--log axes. What slope
do you get, and does it match the claim that iterations grow in proportion
to $\kappa$?

\pe Fit polynomials of degree $1$ through $9$ to T5 using standardised powers,
and plot training cost against degree. Then repeat without standardising and
compare how far each fit gets with a fixed budget of gradient-descent iterations.

\pe Split T5 into eight training points and two held-out points, fit each degree
on the eight, and plot both training and held-out cost against degree on the same
axes. Identify the degree at which the two curves diverge. Keep this plot: it is
the central figure of \cref{ch:ch10}.

\exercisesection{Self-Test}

\qz The condition number of the cost is the ratio of: (a)~the largest to the
smallest feature; (b)~the largest to the smallest Hessian eigenvalue;
(c)~training to test error; (d)~$\alpha$ to $\alpha^\star$.

\qz Dividing each feature by its maximum: (a)~is equivalent to standardisation;
(b)~fixes the ranges but not the centring; (c)~makes the contours circular;
(d)~changes the model's predictions.

\qz Standardising the features of a converged, unregularized linear regression
changes: (a)~the predictions; (b)~the fitted parameter values only;
(c)~neither; (d)~the training data.

\qz After standardisation, the condition number of a two-feature problem is
determined by: (a)~the number of examples; (b)~the correlation between the
features; (c)~the learning rate; (d)~the bias.

\qz Scaling statistics should be computed on: (a)~the whole dataset; (b)~the
training set only; (c)~the test set; (d)~each batch separately.

\qz The model $f(x) = w_1x + w_2x^2 + b$ is: (a)~nonlinear regression;
(b)~linear in the parameters and so fitted by the usual method; (c)~not fittable
by gradient descent; (d)~always better than a straight line.

\qz Adding $x^3$ as a feature when $x$ ranges over $[1,10]$: (a)~has no effect on
conditioning; (b)~adds a feature ranging over three decades and makes scaling
essential; (c)~reduces the condition number; (d)~guarantees a better fit.

\qz A model achieving $J = 0$ on its training set is: (a)~ideal; (b)~evidence
the algorithm worked and nothing else; (c)~impossible; (d)~proof of a bug.

\furtherreading

\textcite{james2021} §3.3 covers the practical extensions of this chapter ---
interaction terms, which are the products of \cref{sec:feature-engineering}, and
polynomial regression --- and §7.1 develops polynomial fits properly, with the
caution about behaviour near the boundary of the data that
\cref{fig:polynomial} illustrates. \textcite{bishop2006} §1.1 is the canonical
treatment of polynomial order and its consequences, and is best read now and
again alongside \cref{ch:ch07}; §3.1 covers the general basis-function view, of
which every transformation in this chapter is an instance.

\bigskip
\noindent\emph{Next:} \Cref{ch:ch06} changes the question. Instead of predicting
a number we predict a category, which breaks the squared-error cost in a way that
no amount of scaling repairs, and requires a new model and a new cost derived
from first principles.
